% Owner: Yogesh Modi. Results: synthesis, resource use and maximum frequency.
% Every number below is taken from synth/reports/ (summary.csv,
% <design>_yosys.txt, <design>_nextpnr.txt, pipeline_breakdown.txt).
\subsection{Synthesis Results}
\label{sec:synth}

\subsubsection{Flow and target}
Both designs were synthesised with the open-source flow of
Section~\ref{sec:litreview}: Yosys~0.71+77 (synthesis) and nextpnr~0.11.1-121
(placement, routing and static timing)~\cite{yogeshWolf2013,yogeshShah2019},
from the OSS CAD Suite build of 22~Nov~2026~\cite{yogeshOSSCAD}. The script
\texttt{synth/run\_synth.sh} reads all files of \texttt{rtl/}, sets the
program with \texttt{chparam -set IMEM\_FILE}, runs \texttt{synth\_ecp5} (or
\texttt{synth\_ice40}) and then nextpnr with the fixed seed~1. Both designs
contain the benchmark \texttt{bench\_sort.hex}. A check in the script confirms
that the 27 program words are in the synthesised instruction ROM. This check
found one problem: the simulation-only loop that fills the instruction memory
with NOPs overwrote the \texttt{\$readmemh} contents in Yosys, so the ROM held
only NOPs. The loop is now excluded from synthesis with
\texttt{`ifndef SYNTHESIS}; simulation is unchanged and all tests still pass.
nextpnr starts at a 50~MHz target. If the target is met it is raised and the
run repeated, so the reported Fmax is the limit of the design: the pipeline
met 50~MHz (56.55~MHz) and failed at 61~MHz (56.55~MHz); the single-cycle
design failed already at 50~MHz (28.43~MHz).

The planned target was the Lattice iCE40 HX8K. The pipeline fits on it (3969
of 7680 logic cells, 16 of 32 block RAMs) and reaches 43.40~MHz. The
single-cycle design does not fit: after synthesis it needs 36\,436 LUTs and
33\,796 flip-flops, and nextpnr stops with 70\,248 required logic cells
(914\,\% of the device). To compare both designs on the same device, the main
results use the larger Lattice ECP5 LFE5U-85F (CABGA381 package, 83\,640
LUTs)~\cite{yogeshECP5DS}, which is supported by the same open flow. The
iCE40 reports are kept in the repository.

\subsubsection{Resources and maximum frequency}
Table~\ref{tab:synth} is \texttt{synth/reports/summary.csv}. LUTs are the
placed logic cells (\texttt{TRELLIS\_COMB}, which include LUTs used as
distributed RAM), FFs are \texttt{TRELLIS\_FF} and BRAMs are 18-kbit
\texttt{DP16KD} blocks.

\begin{table}[!t]
\centering
\caption{Synthesis results on Lattice ECP5 LFE5U-85F (seed 1)}
\label{tab:synth}
\begin{tabular}{lrrrr}
\toprule
Design & LUTs & FFs & BRAMs & Fmax (MHz) \\
\midrule
Pipeline (5-stage) & 3300 & 1358 & 4 & 56.55 \\
Single-cycle       & 7392 & 1040 & 2 & 28.43 \\
\bottomrule
\end{tabular}
\end{table}

The pipeline runs at 1.99 times the clock frequency of the single-cycle
design (56.55 vs.\ 28.43~MHz) and uses fewer logic cells, but 318 more
flip-flops, which are mainly the four pipeline registers.

\subsubsection{Why the memories use (or do not use) block RAM}
The block RAM of the ECP5 and iCE40 has a clocked read
port~\cite{yogeshECP5DS,yogeshICE40DS}, while our memories have a
combinational read. Yosys can still use block RAM if a register directly
follows (or precedes) the read, by merging that register into the RAM. In the
pipeline this is the case for both memories: the IF/ID register follows the
instruction memory and the MEM/WB register follows the data memory, and the
Yosys log reports \emph{merging output FF to cell} for both. Each 1024$\times$32
memory then needs two \texttt{DP16KD} blocks, four in total. In the
single-cycle design the PC register in front of the instruction memory is
merged as its address register (two \texttt{DP16KD}), but the data memory
output goes combinationally into the register file, with no register in
between. On ECP5 it becomes distributed LUT RAM (512 \texttt{TRELLIS\_DPR16X4}
cells), which is why the single-cycle design needs more than twice the logic
cells of the pipeline. The iCE40 has no LUT RAM, so the same memory becomes
32\,768 flip-flops and does not fit. The register file is not mapped to RAM in
either design, mainly because its synchronous reset clears every register at
once, which a RAM cannot do; in the pipeline it is built from 992 flip-flops
(31 registers of 32 bits, \texttt{x0} is constant) and multiplexers. As required, the RTL was not
changed to improve these numbers.

\subsubsection{Per-module breakdown of the pipeline}
Table~\ref{tab:breakdown} is \texttt{synth/reports/pipeline\_breakdown.txt}. It
comes from a separate run with \texttt{synth\_ecp5 -noflatten}, so that every
module keeps its own counts (submodules included). Without flattening, Yosys
cannot merge the pipeline registers into the memories, so here the instruction
memory is a LUT ROM (51 LUTs) and the data memory is LUT RAM; the totals
therefore differ from Table~\ref{tab:synth}.

\begin{table}[!t]
\centering
\caption{Pipeline resources per module (ECP5, not flattened)}
\label{tab:breakdown}
\begin{tabular}{llrrrr}
\toprule
Stage & Module & LUT4 & FF & Carry & LUT RAM \\
\midrule
IF  & \texttt{pc\_reg}          &    0 &   32 &  0 &   0 \\
IF  & \texttt{instr\_mem}       &   51 &    0 &  0 &   0 \\
ID  & \texttt{reg\_file}        & 2872 &  992 &  0 &   0 \\
ID  & \texttt{imm\_gen}         &  123 &    0 &  0 &   0 \\
ID  & \texttt{control\_unit}    &   41 &    0 &  0 &   0 \\
EX  & \texttt{ex\_stage}        & 1348 &    0 & 95 &   0 \\
MEM & \texttt{data\_mem}        & 3107 &    0 &  0 & 512 \\
--  & \texttt{if\_id\_reg}      &    1 &   64 &  0 &   0 \\
--  & \texttt{id\_ex\_reg}      &    1 &  158 &  0 &   0 \\
--  & \texttt{ex\_mem\_reg}     &    0 &   73 &  0 &   0 \\
--  & \texttt{mem\_wb\_reg}     &    0 &   71 &  0 &   0 \\
--  & \texttt{forwarding\_unit} &   22 &    0 &  0 &   0 \\
--  & \texttt{hazard\_unit}     &   22 &    0 &  0 &   0 \\
\midrule
\multicolumn{2}{l}{Whole design} & 7652 & 1390 & 110 & 512 \\
\bottomrule
\end{tabular}
\end{table}

Apart from the memories, the register file is the largest block: its two
32-to-1 read multiplexers of 32 bits and the write-first bypass take 2872 LUTs.
The EX stage (ALU with shifter, branch comparator and the forwarding
multiplexers) takes 1348 LUTs and almost all carry cells. The four pipeline
registers add 366 flip-flops, and the forwarding and hazard units are small
(22 LUTs each), so hazard handling costs little hardware.

\subsubsection{Critical path}
In the pipeline (ECP5) the critical path is 17.68~ns (10.30~ns logic,
7.38~ns routing). It starts at the output register of the data-memory block
RAM, which is the MEM/WB load-data register, goes through the write-back
multiplexer, the MEM/WB forwarding path into the ALU, and ends at the EX/MEM
register. The limiting stage is therefore EX, lengthened by the forwarded
load value from WB. On the iCE40 HX8K the critical path (23.04~ns) also starts
at the data-memory read register, but ends at the flush input of the ID/EX
register: the forwarded value decides a branch in EX, and the taken-branch
signal flushes the ID/EX register. In the single-cycle design the critical
path is 35.17~ns (9.81~ns logic, 25.36~ns routing): from the instruction ROM
output, through the register file read, the ALU and the distributed-RAM read of
the data memory, to the register-file write, i.e.\ the whole datapath in one
cycle. Its large routing share is most likely due to the 512 LUT-RAM cells
of the data memory, which are spread over a large area of the device.

\subsubsection{Interpretation}
Pipelining roughly halves the clock period (17.68 vs.\ 35.17~ns), not divides
it by five, because the stages are not balanced: EX with forwarding is much
longer than IF or WB, and the pipeline registers add clock-to-output and setup
time~\cite{yogeshHennessyQA,yogeshHartstein2002}. The pipeline is also smaller
on this FPGA, which is not true in general: here it comes from the pipeline
registers that let both memories use block RAM. The execution time of the
benchmarks, which combines these Fmax values with the measured CPI, is given
in the performance results.
